Hamming Quasi-Cyclic~\cite{melchor2018hamming} is a fourth round candidate in the NIST Post-Quantum cryptography standardization process.
HQC is a code-based cryptosystem, that includes comprises both a public-key encryption (\textsf{HQC.PKE}) and a
key encapsulation mechanism (\textsf{HQC.KEM}).\\
In this section we will describe the inner workings of HQC and show its theoretical security; moreover,
we will give an overview of the currently known attacks against HQC and their practical impact.

\section{Encryption scheme}\label{sec:encryption}
We can think of a generic encryption scheme as a tuple of four polynomial time algorithms:
\begin{itemize}
    \item \textsf{Setup}$(1^\lambda)$: generates the global parameters \textsf{param} based on the security parameter $\lambda$.
    \item \textsf{Keygen(param)}: generates a public key \textsf{pk} and a secret key \textsf{sk} based on the global parameters \textsf{param}.
    \item \textsf{Encrypt(pk, \textbf{m}, $\theta$)}: outputs the ciphertext \textbf{c} for the message \textbf{m} with the randomness $\theta$ under the public key \textsf{pk}.
    \item \textsf{Decrypt(sk, \textbf{c})}: outputs the message \textbf{m} for the ciphertext \textbf{c} (or $\bot$ if the ciphertext is invalid).
\end{itemize}

TODO: add security proof?

\section{Public Key Encryption}\label{sec:pke}
HQC uses two different codes:
\begin{itemize}
    \item A publicly available decodable code $\mathcal{C} = [n, k]$ generated by $\mathbf{G} \in \mathds{F}_2^{k \times n}$  with an efficient decoding algorithm 
    $\mathcal{C}.\textsf{Decode}()$ which can correct up to $\delta$ errors.
    \item A random double-circulant $[2n,n]$ code with parity-check matrix $\mathbf{H} = (\mathbf{I_n, circ(h)})$. 
\end{itemize}

The \textsf{HQC.PKE} defines the four algorithms of an encryption scheme as follows:
\begin{itemize}
    \item \textsf{Setup}$(1^\lambda)$: generates the set of parameters \textsf{param} = $(n, k, \delta, \omega, \omega_r, \omega_e)$
    \item \textsf{Keygen(param)}: samples $\mathbf{h} \xleftarrow{\mathdollar} \mathcal{R}$ and 
    $(\mathbf{x}, \mathbf{y}) \xleftarrow{\mathdollar} \mathcal{R} \times \mathcal{R}$ such that $\texttt{wt}(\mathbf{x}) = \texttt{wt}(\mathbf{y}) = \omega$;
    the keypair is defined as \textsf{sk} = $(\mathbf{x}, \mathbf{y})$ and \textsf{pk} = $(\mathbf{h}, \mathbf{s} = \mathbf{x} + \mathbf{h\cdot y})$.
    \item \textsf{Encrypt(pk, \textbf{m}, $\theta$)}: samples $\mathbf{e} \xleftarrow{\mathdollar} \mathcal{R}$, $(\mathbf{r_1}, \mathbf{r_2}) \xleftarrow{\mathdollar} \mathcal{R} \times \mathcal{R}$ such that
     $\texttt{wt}(\mathbf{e}) = \omega_e$ and $\texttt{wt}(\mathbf{r_1}) = \texttt{wt}(\mathbf{r_2}) = \omega_e$;  the ciphertext is defined as the pair
     $(\mathbf{u}, \mathbf{v})$, where $\mathbf{u} = \mathbf{r_1} + \mathbf{h}\cdot\mathbf{r_2}$ and $\mathbf{v} = \mathbf{mG} + \mathbf{s\cdot r_2} + \mathbf{e}$.
    \item \textsf{Decrypt(sk, \textbf{c})}: outputs $\mathcal{C}.\textsf{Decode}(\mathbf{v} - \mathbf{u\cdot y})$.
\end{itemize}

The idea at the core of \textsf{HQC.PKE} is to encode the message ($\mathbf{mG}$) and then add random errors that can be removed only by knowing the secret key.\\

TODO: add security proof?

\section{Key Encapsulation Mechanism}\label{sec:kem}
In practice, public key encryption is not used to encrypt large amounts of data because of its high exectution time compared to symmetric encryption.
For this reason, we use public key encryption to exchange a shared secret key, which is then used to encrypt the actual data: the practice
of using a public key encryption scheme to exchange a shared secret key is referred to as key encapsulation mechanism (KEM).\\
The Fujisaki-Okamoto transformation~\cite{fujisaki2013secure} is a technique that allows to transform a \textit{weak} IND-CPA encryption scheme
into a \textit{stronger} IND-CCA2 key encapsulation mechanism.

Let $\mathcal{E}$ an instance of \textsf{HQC.PKE} and $\mathcal{G, H}$ and $\mathcal{K}$ hash functions. The \textsf{HQC.KEM} scheme is defined as follows:
\begin{itemize}
    \item \textsf{Setup}$(1^\lambda)$: generates the set of parameters as in \textsf{HQC.PKE.Setup}, where $k$ will be the key size.
    \item \textsf{Keygen(param)}: generates a keypair as in \textsf{HQC.PKE.Keygen}.
    \item \textsf{Encapsulate(pk)}: sample $\mathbf{m} \xleftarrow{\mathdollar} \mathds{F}_2^k$ to derive the randomness $\theta = \mathcal{G}(\mathbf{m})$.
    Generate the ciphertext as $\mathbf{c} = (\mathbf{u, v}) \leftarrow \mathcal{E}.\textsf{Encrypt}(\textsf{pk}, \mathbf{m}, \theta)$ and derive the symmetric key
    as $K \leftarrow \mathcal{K}(\mathbf{m, c})$; send $(\mathbf{c, d} = \mathcal{H}(\mathbf{m})).$
    \item \textsf{Decapsulate}(\textsf{sk}, $\mathbf{c, d}$): decrypt the ciphertext as $\mathbf{m^\prime} = \mathcal{E}.\textsf{Decrypt}(\textsf{sk}, \mathbf{c})$, 
    compute $\theta^\prime = \mathcal{G}(\mathbf{m}^\prime)$ and re-encrypt to get $\mathbf{c^\prime} \leftarrow \mathcal{E}.\textsf{Encrypt}(\textsf{pk}, \mathbf{m^\prime}, \theta)$.
    If $\mathbf{c \neq c^\prime}$ or $\mathbf{d} \neq \mathcal{H}(\mathbf{m^\prime})$ return $\bot$, otherwise derive the shared key $\mathcal{K}(\mathbf{m^\prime, c^\prime})$.
\end{itemize}

TODO: what happens if we do not check $\mathbf{d} == \mathcal{H}(\mathbf{m^\prime})$?
